Каскадные системы позволяют адаптивно распределять вычислительные ресурсы, передавая часть запросов от дешёвого предиктора более дорогому. Эффективность такой системы зависит от способности маршрутизатора оценить ожидаемое изменение качества ответа: дорогая модель может как исправить ошибку, так и заменить правильный ответ неправильным. Существующие маршрутизаторы используют различные обучающие цели и параметризации, однако для выбора того или иного подхода для новой задачи важно учитывать доступные данные, структуру ошибок моделей и ограничения эксплуатации. Цель исследования заключается в выявлении условий, при которых различные способы обучения маршрутизатора будут обеспечивать наилучшее соотношение качества и вычислительных затрат. Основная задача - разработать систему рекомендаций по выбору маршрутизатора на основе характеристик данных и используемых моделей. В работе будет рассматривается каскад из двух фиксированных классификаторов и маршрутизатора, получающего на вход признаки объекта и результат дешёвой модели. Для сравнения будут выбраны такие политики маршрутизации как прямая регрессия прироста качества, разность оценок правильности моделей, оценивание вероятностей исправления и ухудшения качества ответа, неуверенность только дешевой модели, сложность признаков входного объекта. Переносимость сформированных рекомендаций планируется проверить на новых наборах данных и парах классификаторов, не использованных при выявлении закономерностей и выборе параметров методов. Итогом работы должны стать экспериментально обоснованные правила выбора способа обучения маршрутизатора, позволяющие определить, когда дополнительная сложность оправданна, а когда сопоставимое качество достигается более простым решением.
